% Part: computability
% Chapter: recursive-functions
% Section: composition

\documentclass[../../../include/open-logic-section]{subfiles}

\begin{document}

\olfileid{cmp}{rec}{com}
\olsection{ترکيب}

که~$f$ او~$g$ د طبيعي عددونو دوه يوځاييزې تابعې وي، نو ترکيب ئې
کولے شو: $h(x) = f(g(x))$۔ په دې حالت کښې نوې تابع~$h(x)$ له
تابعو~$f$ او~$g$ څخه په \emph{ترکيب} تعريف شوې ده۔ غواړو دا د يو
څخه د زياتو آرګومېنټونو تابعو ته عمومي کړو۔

د دې يوه لاره داسې ده: فرض کړئ~$f$ يوه~$k$-ځاييزه تابع ده، او
$g_0$، \dots،~$g_{k-1}$ هغه~$k$ تابعې دي چې ټولې~$n$-ځاييزې دي۔
بيا يوه نوې~$n$-ځاييزه تابع~$h$ داسې تعريفولے شو:
\[
h(x_0, \dots, x_{n-1}) =
f(g_0(x_0, \dots, x_{n-1}), \dots, g_{k-1}(x_0, \dots, x_{n-1}))
\]
که~$f$ او ټولې~$g_i$ محاسبه کېدونکې وي، نو~$h$ هم محاسبه کېدونکې
ده: د~$h(x_0, \dots, x_{n-1})$ د محاسبې دپاره لومړے د هر
$i = 0$، \dots،~$k-1$ دپاره ارزښتونه~$y_i = g_i(x_0, \dots,
x_{n-1})$ محاسبه کړئ۔ بيا دا ارزښتونه~$f$ ته ورکړئ، څو
$h(x_0, \dots, x_{n-1}) = f(y_0, \dots, y_{k-1})$ محاسبه کړئ۔
\textbf{د ژباړې د نسخې سمون (OLCMP-003):} سرچينې په وروستۍ محاسبې
کښې د تابع د آرګومېنټونو لړۍ د تابعو د شمېر له مخې پای ته رسولې وه؛
دلته لړۍ د تابع د ځاييزې کچې له مخې بشپړه شوې ده۔

ښايي دا د دې خبرې له حده زيات محدود بيان ښکاره شي چې له ځينو شته
تابعو څخه د نوې تابع د محاسبې پر مهال څۀ پېښېږي۔ يوه خبره دا ده چې
کله ناکله د يوې تابع ټول آرګومېنټونه نۀ کاروو؛ لکه هغه وخت چې د
$\Add$ په بنسټيز بازګشتي تعريف کښې مو د کارولو دپاره
$g(x, y, z) = \Succ(z)$ تعريف کړه۔ فرض کړئ د لاندې تابعو د کارولو
اجازه لرو:
\[
\Proj{n}{i}(x_0, \dots, x_{n-1}) = x_i
\]
تابعو~$\Proj{n}{i}$ ته \emph{پروجکشن} تابعې وايي:
$\Proj{n}{i}$ يوه~$n$-ځاييزه تابع ده۔ بيا~$g$ داسې تعريفېدے شي:
\[
g(x, y, z) = \Succ(\Proj{3}{2}(x, y, z)).
\]
دلته د~$f$ ونډه~$1$-ځاييزه تابع~$\Succ$ ترسره کوي، نو~$k=1$ دے۔ او
يوه~$3$-ځاييزه تابع~$\Proj{3}{2}$ لرو چې د~$g_0$ ونډه ترسره کوي۔
پايله يوه~$3$-ځاييزه تابع ده چې د درېيم آرګومېنټ راتلونکی ورکوي۔
\textbf{د ژباړې د نسخې سمون (OLCMP-004):} سرچينې د پروجکشن تابعو
له تعريف سملاسي وروسته په نوم کښې د تابعو شمېر د ځاييزې کچې پر ځای
کارولی و؛ دلته نوم له هماغې تعريف شوې ځاييزې کچې سره سم شوے دے۔

پروجکشن تابعې دا اجازه هم راکوي چې د آرګومېنټونو په بيا ترتيبولو يا
سره يو کولو نوې تابعې تعريف کړو۔ د بېلګې په توګه، تابع
$h(x) = \Add(x, x)$ داسې تعريفېدے شي:
\[
h(x_0) = \Add(\Proj{1}{0}(x_0),\Proj{1}{0}(x_0)).
\]
دلته~$k=2$، $n=1$، د~$f(y_0,y_1)$ ونډه~$\Add$ ترسره کوي، او د
$g_0(x_0)$ او~$g_1(x_0)$ دواړو ونډې~$\Proj{1}{0}(x_0)$، يعنې
يوځاييزه پروجکشن تابع (يا عينيت تابع)، ترسره کوي۔

که~$f(y_0, y_1)$ هغه تابع وي چې له مخکښې ئې لرو، نو تابع
$h(x_0, x_1) = f(x_1, x_0)$ داسې تعريفولے شو:
\[
h(x_0, x_1) = f(\Proj{2}{1}(x_0, x_1),\Proj{2}{0}(x_0, x_1)).
\]
دلته~$k=2$، $n = 2$، او د~$g_0$ او~$g_1$ ونډې په ترتيب سره
$\Proj{2}{1}$ او~$\Proj{2}{0}$ ترسره کوي۔

ښايي دا اندېښنه هم ولرئ چې د ټولو~$g_0$، \dots،~$g_{k-1}$ ځاييزه
کچه بايد يو شان~$n$ وي۔ (په ياد ولرئ چې د يوې تابع
\emph{ځاييزه کچه} د آرګومېنټونو شمېر دے؛ د~$n$-ځاييزې تابع ځاييزه
کچه~$n$ ده۔) خو د پروجکشن تابعو زياتول اړينه انعطاف‌پذيري برابروي۔
د بېلګې په توګه، فرض کړئ~$f$ او~$g$، $3$-ځاييزې تابعې دي او~$h$
هغه~$2$-ځاييزه تابع ده چې داسې تعريف شوې ده:
\[
h(x,y) = f(x,g(x,x,y),y).
\]
د~$h$ تعريف له پروجکشن تابعو سره داسې بيا ليکل کېدے شي:
\[
h(x,y) = f(\Proj{2}{0}(x,y), g(\Proj{2}{0}(x,y), \Proj{2}{0}(x,y),
\Proj{2}{1}(x,y)), \Proj{2}{1}(x,y)).
\]
بيا~$h$ له~$\Proj{2}{0}$، $l$ او~$\Proj{2}{1}$ سره د~$f$ ترکيب
دے، چې
\[
l(x,y) = g(\Proj{2}{0}(x,y),\Proj{2}{0}(x,y),\Proj{2}{1}(x,y)),
\]
يعنې~$l$ له~$\Proj{2}{0}$، $\Proj{2}{0}$ او~$\Proj{2}{1}$ سره د
$g$ ترکيب دے۔

\end{document}
